\section{Xây dựng giao diện người dùng (Frontend --- React + Vite)}
\label{sec:ch3}

\subsection{Cấu trúc thư mục \code{frontend/}}

Dự án frontend được khởi tạo bằng Vite với mẫu React. Các tệp do người lập trình viết hoặc chỉnh sửa
nằm trong \code{src/}:

\begin{lstlisting}[style=ShellStyle, caption={Cấu trúc thư mục frontend}]
frontend/
|-- index.html               # root HTML page, loads /src/main.jsx
|-- package.json             # dependencies and scripts (dev, build, lint)
|-- vite.config.js           # Vite config (React plugin)
|-- public/                  # favicon and static icons
`-- src/
    |-- main.jsx             # entry point: renders <App /> into #root
    |-- App.jsx              # main component: UI + business logic
    |-- App.css              # application styles
    |-- index.css            # global base styles
    |-- config/
    |   `-- contract.js      # CONTRACT_ADDRESS, CONTRACT_ABI, TOUR_DEFAULT
    `-- utils/
        `-- web3Client.js    # MetaMask / web3.js helper functions
\end{lstlisting}

Việc tách tệp tuân theo nguyên tắc phân tầng: \code{config/contract.js} chứa \emph{dữ liệu cấu hình} (địa chỉ và ABI
của hợp đồng), \code{utils/web3Client.js} chứa \emph{lớp truy cập blockchain}, còn \code{App.jsx} chỉ lo trạng thái và giao diện.
Khi triển khai lại hợp đồng, chỉ cần đổi \code{CONTRACT\_ADDRESS} trong một tệp.

\subsection{Thiết kế giao diện (\code{App.jsx}, \code{App.css})}

Giao diện là một trang đơn (single page) gồm các khối từ trên xuống:

\begin{enumerate}[leftmargin=1.8em]
    \item \textbf{Thanh đầu trang (topbar):} tên ứng dụng và khu vực ví. Khi chưa kết nối hiện nút \emph{Kết nối MetaMask};
          khi đã kết nối hiện ``viên thuốc'' (wallet pill) gồm vai trò (Điều hành/Khách hàng), tên mạng, số dư ETH và địa chỉ rút gọn.
    \item \textbf{Biểu ngữ chế độ:} cho biết đang ở \emph{chế độ On-Chain} (đã khai báo \code{CONTRACT\_ADDRESS}) hay chế độ mô phỏng.
    \item \textbf{Thông báo trạng thái (alert):} xanh = thành công, đỏ = lỗi, xanh dương = thông tin; có nút đóng.
    \item \textbf{Ba thẻ thống kê:} giá vé/người, chỗ trống còn lại, và số vé đã đặt (xem Mục~\ref{sec:stats}).
    \item \textbf{Bảng thao tác theo vai trò:} form tạo tour (Điều hành) hoặc form đặt chỗ (Khách hàng).
\end{enumerate}

\code{App.css} dùng CSS thuần với bố cục \code{flex}/\code{grid} (ví dụ \code{.stats-grid} dùng
\code{repeat(auto-fit, minmax(220px, 1fr))} để tự co giãn theo độ rộng màn hình), các lớp màu cho thông báo
(\code{.alert-success}, \code{.alert-error}, \code{.alert-info}) và nút (\code{.btn-primary}, \code{.btn-secondary}).

\subsection{Quản lý trạng thái}

\code{App.jsx} dùng hook \code{useState} để giữ các trạng thái chính:

\begin{table}[H]
\centering
\caption{Các trạng thái chính của component \code{App}}
\begin{tabularx}{\linewidth}{|l|X|}
\hline
\rowcolor{navyblue}\textcolor{white}{\textbf{Trạng thái}} & \textcolor{white}{\textbf{Ý nghĩa}} \\
\hline
\code{account}, \code{balanceEth}, \code{chainId} & Ví đang kết nối, số dư ETH, mạng hiện tại. \\
\hline
\rowcolor{rowgray}
\code{operatorAddress} & Địa chỉ điều hành, đọc từ \code{operator()} của hợp đồng. \\
\hline
\code{totalSeats}, \code{availableSeats}, \code{ticketPriceEth} & Thông tin tour đọc từ hợp đồng. \\
\hline
\rowcolor{rowgray}
\code{myBookings} & Số vé của ví đang kết nối (từ \code{getMyBookings()}). \\
\hline
\code{seatsToBook}, \code{newTourSeats}, \code{newTourPrice} & Giá trị người dùng đang nhập trong hai form. \\
\hline
\rowcolor{rowgray}
\code{loading}, \code{statusMessage} & Trạng thái chờ giao dịch và thông báo hiển thị cho người dùng. \\
\hline
\end{tabularx}
\end{table}

\subsection{Hiển thị giá tour, chỗ trống và vé đã đặt}
\label{sec:stats}

Ba thẻ thống kê (class \code{stat-card}) đáp ứng ràng buộc số 4 của đề:

\begin{itemize}[leftmargin=1.5em]
    \item \textbf{Giá vé/người:} \code{ticketPriceEth} --- đọc \code{ticketPrice()} (đơn vị wei) rồi đổi sang ETH bằng
          \code{web3.utils.fromWei}.
    \item \textbf{Chỗ trống còn lại:} \code{availableSeats / totalSeats} --- dữ liệu từ \code{getAvailableSeats()} và \code{totalSeats()}.
    \item \textbf{Vé đã đặt:} với khách hàng là \code{myBookings} (từ \code{getMyBookings()}); với điều hành, thẻ này
          đổi nhãn thành \emph{Số chỗ đã được đặt} và hiển thị tổng $\texttt{totalSeats} - \texttt{availableSeats}$ của toàn tour.
\end{itemize}

\lstinputlisting[style=JSStyle, firstline=290, lastline=316,
  caption={Ba thẻ thống kê trong \code{App.jsx}}]{Code/App.jsx}

\subsection{Cơ chế tự động nhận diện vai trò (\code{isOperator})}

Thay vì bắt người dùng chọn vai trò, giao diện tự suy ra vai trò từ ví đang kết nối:

\lstinputlisting[style=JSStyle, firstline=40, lastline=43,
  caption={Xác định vai trò bằng so sánh địa chỉ ví}]{Code/App.jsx}

Biến \code{isOperator} là \code{true} khi địa chỉ ví kết nối trùng với \code{operator()} lấy từ hợp đồng
(so sánh sau khi chuyển về chữ thường, vì địa chỉ Ethereum có thể viết hoa/thường khác nhau do chuẩn checksum).
Dựa vào đó, phần hiển thị chia thành ba nhánh (\code{App.jsx}, khối \code{role-panel-wrap}):

\begin{table}[H]
\centering
\caption{Ba trạng thái giao diện theo vai trò}
\begin{tabularx}{\linewidth}{|l|l|X|}
\hline
\rowcolor{navyblue}\textcolor{white}{\textbf{Điều kiện}} & \textcolor{white}{\textbf{Giao diện}} & \textcolor{white}{\textbf{Nội dung}} \\
\hline
\code{!account} & Chưa kết nối & Hướng dẫn và nút \emph{Kết nối MetaMask ngay}. \\
\hline
\rowcolor{rowgray}
\code{isOperator} & Điều hành & Form \emph{Mở / Tạo Tour} (tổng số chỗ, giá vé) gọi \code{createTour()}. \\
\hline
còn lại & Khách hàng & Form \emph{Đặt chỗ Tour} (số chỗ, tổng thanh toán dự kiến) gọi \code{bookSeats()}. \\
\hline
\end{tabularx}
\end{table}

\shot{create_tour.png}{Giao diện Điều hành: ví Account 3 được nhận diện là \emph{Điều hành}, hiển thị form tạo tour (tổng số chỗ, giá vé/người)}{0.95}

\shot{cus_ui.png}{Giao diện Khách hàng: ví Account 2 được nhận diện là \emph{Khách hàng}, hiển thị form đặt chỗ và tổng thanh toán dự kiến}{0.95}

Khi người dùng đổi tài khoản trong MetaMask, sự kiện \code{accountsChanged} được xử lý lại (\ref{sec:ch4}),
\code{isOperator} được tính lại và giao diện tự chuyển giữa hai chế độ mà không cần tải lại trang.

\subsection{Kiểm tra dữ liệu đầu vào ở giao diện}

Trước khi gửi giao dịch, các hàm \code{handleCreateTour} và \code{handleBookSeats} kiểm tra: đã kết nối ví chưa;
số chỗ là số nguyên dương; số chỗ đặt không vượt \code{availableSeats}; giá vé lớn hơn 0; và (với tạo tour) ví hiện tại là điều hành.
Ô nhập cũng đặt \code{min} và \code{max} (\code{max=availableSeats}) để trình duyệt chặn giá trị sai ngay khi nhập.
Các kiểm tra này giúp báo lỗi sớm và tránh tốn phí gas cho giao dịch chắc chắn thất bại; tuy vậy,
\textbf{lớp bảo vệ cuối cùng vẫn là các câu lệnh \code{require} trong smart contract} (\ref{sec:ch5}).

\subsection{Chế độ mô phỏng}

Nếu \code{CONTRACT\_ADDRESS} rỗng, ứng dụng chạy ở \emph{chế độ mô phỏng}: dữ liệu giữ trong bộ nhớ (\code{mockBookingsByWallet})
và không gửi giao dịch lên chuỗi. Chế độ này giúp thiết kế giao diện khi chưa có hợp đồng; bản nộp dùng chế độ \emph{On-Chain}
(thanh biểu ngữ màu xanh trong các ảnh chụp).
